% ===== §1 The Phenomenon and the Paradox =====
\section{The Phenomenon and the Paradox}
\label{sec:phenomenon}

\noindent
Consider two people who love each other, and one of whom is, in the ordinary worldly sense, the stronger---the quicker to act, the more fluent in plans, the one to whom solutions occur first and arrangements come easily. He loves her, and his love takes the form his strength makes natural: he provides, he foresees, he smooths the road before she reaches the rough part of it, he settles the difficulty before she has finished framing it. Nothing he does is to her detriment in any register the usual theories can name. He takes nothing from her; he deceives her in nothing; he honours her, listens to her, would account to her for anything she asked. By every available measure the relation prospers, and she, asked, would say she is happy and is not lying.

And yet something is wrong, and the wrongness is hers. Somewhere in the smooth provision a space has closed. The decisions that were to be made were made before she came to them; the meanings that were to be found were named before she found them; the road was paved before she could choose its direction. She is not oppressed---there is no oppressor---and she is not deprived, for she wants for nothing the relation can supply. She is something for which the language of justice has no ready word: she has been generated \emph{for}, so thoroughly and so well, that there is nothing left for her to generate, and a self that lived by generating has quietly gone still. This paper is written to give that stillness a name and a theory.

\subsection{The datum}
\label{sub:phen-datum}

The phenomenon, stated plainly, is this. In a shared life understood as a field of co-generation---a coupling out of which value arises between two people rather than being brought to it ready-made by either---one party may generate so far ahead of the other, so fast and so completely, that the other finds no remaining occasion for her own generation. The harm is not that the fast-generating party produces \emph{badly}, nor that he produces \emph{for himself}; he produces well, and for the relation, and often most of all for her. The harm is that the space of generation is shared and finite, and that a sufficiently rapid and total generation by one pole leaves the other pole's generativity with nowhere to act. What withers is not the other's holdings, which may grow, nor her standing, which may be high, but her activity---the exercise, now without occasion, of her own capacity to make value at her own pace.

It is essential to the datum that the withering is not chosen as a deprivation and not experienced, at first, as one. The space does not close against resistance; it closes \emph{around} a gratitude. Each particular act of provision is a kindness, received as a kindness; it is only their sum, over time, that amounts to an occupation. This is why the harm is so often invisible to both parties at once: the stronger pole sees only a sequence of devotions, each defensible, each loving; the weaker pole feels only a slow and unaccountable diminishment of her own activity, with no event to point to and no one to blame. There is no scene of the crime, because there is no crime---only a good, done too completely, in a space that could not hold all of it without displacing her.

\subsection{The paradox that defeats existing theory}
\label{sub:phen-paradox}

The datum is built so as to defeat, one by one, every criterion the theory of justice ordinarily brings to a relation. This is not an accident of the example but the point of it: the phenomenon is interesting precisely because it falls into the blind spot common to all the standard diagnoses.

There is no \emph{appropriation}. The stronger pole takes nothing from the weaker; on the contrary, he gives. The vocabulary of exploitation, of extraction, of one party living on another's unrecognised labour, finds nothing to grip: the flow of value runs toward the effaced party, not away from her.

There is no \emph{domination}. The stronger pole issues no command, brooks no resistance because he meets none, imposes nothing she refuses. He serves. The vocabulary of coercion and subjection, of a will bent to another's, finds nothing to grip either: there is no contest of wills in which one prevails, only a devotion that forecloses by filling rather than by forcing.

There is no \emph{failure to recognise or to hear}. The stronger pole sees her, values her, attends to her, would credit her every contribution. The vocabulary of misrecognition and of epistemic injustice finds nothing to grip: she is not unseen, not unheard, not denied her standing as a knower. She is acknowledged in full---and effaced in full---at once.

\begin{claim}[The paradox of devotion]
Generative effacement is a harm that arrives through devotion and survives the absence of every standard wrong. There is no appropriation, for the stronger pole gives rather than takes; no domination, for he serves rather than commands; no failure of recognition or of hearing, for he honours and attends. Every criterion of injustice built upon possession, coercion, or non-acknowledgement reports that nothing is wrong---and the harm is nonetheless real, and is the weaker pole's. A theory adequate to it must therefore find the wrong at a level beneath possession, coercion, and acknowledgement altogether.
\end{claim}

\subsection{Subordination completed in sincere endorsement}
\label{sub:phen-sincere}

One move remains, and it is the one that closes every exit. It might be supposed that the harm, though it survives the absence of the standard wrongs, would at least announce itself in the effaced party's own dissent---that she, whatever the outside observer misses, would feel and could voice her diminishment, so that her objection might serve as the criterion her circumstances otherwise lack. The datum denies even this. The effaced party frequently does not object; she endorses. She is grateful for the provision, attached to the provider, and would, asked whether she would have it otherwise, sincerely say no. Her subordination is completed not in false consciousness imposed from without but in an endorsement that is, by every test available to her, authentically her own.

This feature---that the harm survives even the sincere endorsement of the one wronged---sets the problem the rest of the paper must solve. A wrong that survives not only the absence of the standard wrongs but the \emph{sincere endorsement of the one wronged} cannot be located in anything the will refuses, for her will does not refuse it---it embraces it. The harm must lie somewhere the will does not reach: at the level not of what she chooses but of what she \emph{is}, not of her endorsement but of her being. To show that there is such a level, that an injustice can be done there, and that it can be done even by love and even with consent---and then to say what justice, at that level, would require---is the work this paper undertakes. It begins, in the next section, by locating the level: by showing where, beneath the injustices the theory of justice already knows, the ground of being lies.
